\section{Related work}
\label{sec:related}

This paper sits in a literature that has established most of its framing. We group that work by
theme, say what each strand did, and end with what is left for us.

\subsection{Fiscal risk, contingent liabilities and public credit}

The tradition this paper belongs to treats a government's exposure as a balance sheet rather than
a deficit. \citet{Polackova1998} set out the fiscal risk matrix that distinguishes explicit from
implicit and direct from contingent obligations, and argued that the obligations outside the
budget are the ones that destabilize it. \citet{Cebotari2008} developed the measurement and
disclosure practice that followed, and \citet{BovaRuizArranzToscaniTure2019} showed empirically
that realizations of contingent liabilities are large, frequent and concentrated in the financial
sector. The valuation side is due to \citet{Lucas2012}, who set out how government policies and
projects should be priced when the risk is not diversifiable, to \citet{LucasMcDonald2010}, who
valued the guarantees behind the housing agencies, and to \citet{Lucas2016}, who showed that
federal credit programs are fiscal policy conducted through a balance sheet whose cost is not
what the budget records.

What that tradition does not have is a systematic decomposition of claims by the \emph{income}
that services them. A fiscal risk matrix classifies an obligation by whether the government must
pay it and under what trigger; labor backing accounts ask what income stream has to keep flowing
for it to be serviced, and who is left holding it if that stream stops. The accounts sit beside
the matrix rather than replacing it, adding the servicing dimension and applying it to the whole
claim structure rather than to the government's own obligations alone, which is what makes the
two sides comparable.

A related strand prices the link in the other direction. \citet{AcharyaDrechslerSchnabl2014}
showed that bank rescues transfer risk to the sovereign, \citet{GennaioliMartinRossi2014} that
domestic banks holding sovereign debt transmit default back into credit, and
\citet{FarhiTirole2018} that the two feed each other; \citet{DellAricciaEtAl2018} survey the
policy problem. That literature connects the sovereign to banks through bank holdings of
government debt. The connection this paper measures runs the other way and through a third party:
both the sovereign and the banks are exposed to the same household income stream, and neither
exposure appears on the other's balance sheet.

\subsection{Financial accounts and who-to-whom structure}

The data tradition begins with \citet{Copeland1952}, whose money flows study established that a
financial system can be described as a matrix of claims between sectors. The Federal Reserve's
financial accounts \citep{FRBZ1} are the modern form of it, and the counterparty structure has
been used to build sector networks \citep{CastrenRancan2014} and, distributionally, to place
household wealth inside the aggregate \citep{BattyEtAl2019}. The Federal Reserve's own debt
service ratio is the closest published statistic to our object: it measures household debt
service against all household income.

Our accounts extend that structure in one way. Who-to-whom accounting records the counterparties
of a claim but not which income stream services it; we add that dimension across thirteen classes
and cross it with the holder structure the accounts already carry. The debt service ratio does
the servicing calculation for households alone and against all of their income, while we do it
for the whole claim structure and against labor income, which is what makes the government's
position visible at all.

\subsection{Household balance sheets and the distribution of debt}

The household side of our measurement draws on a literature that has established how income
shocks, liquid buffers and distress travel together
\citep{MianRaoSufi2013,GanongNoel2020,BhuttaBrickerDettlingKelliherLaufer2019,MerikullRoom2020}.
\citet{BartscherKuhnSchularickSteins2020} is the nearest work in spirit, tracing the joint
history of inequality, household debt and financial fragility in the United States since 1950.
Their object is the household balance sheet over time; ours is the claim stock seen from the
other side, which lets the same debt be attributed to the holder who carries it rather than only
to the household that owes it.

That difference matters for what a distributional statement can say. Placing household wealth
inside the aggregate \citep{BattyEtAl2019} answers how much each group owns. Splitting the
wage-backed claim stock by wage quintile answers something adjacent and separate: how much of the
income that services the system each group has to supply. The two readings diverge, and the
divergence is the result we report in Section~\ref{sec:households}.

\subsection{What is ours}

\begin{quote}
The fiscal mechanism is established, and this paper does not claim it. What this paper adds is
the measured holder structure of the claims that wages pay, on both sides and on the same basis;
a series for that structure running back to 1947; the distribution of the wage-backed stock
across the wage distribution that services it; and a public claims register in which every
reported quantity carries its status.
\end{quote}

\paragraph{What a precedent would have to do, and what the nearest work does instead.} We fixed
in advance what would count as prior construction of this object: a decomposition that takes a
national claim stock, classifies each claim by the income that services it, crosses that with the
holder structure, and reports an aggregate share. A search under that protocol, with the kill
criterion stated before it was run, located none. The nearest work and the condition each fails:

\begin{description}[leftmargin=2.2em,style=nextline,itemsep=2pt]
  \item[Debt service ratios] measure a flow of required payments against income, for households,
  against disposable income. Ours is a stock decomposition, across all sectors, against labor
  income. No servicing-income classification of the claim stock.
  \item[Distributional national accounts] \citep{BattyEtAl2019} distribute wealth and impute debt
  across the income distribution, answering who owns and owes rather than which income services
  what, and they are not crossed with the holder structure.
  \item[Fiscal risk matrices] \citep{Polackova1998,Cebotari2008} classify an obligation by its
  trigger, and cover the government's own obligations rather than the whole claim structure.
  \item[Who-to-whom accounting] \citep{Copeland1952,FRBZ1} classifies by instrument and
  counterparty. This is the structure we extend; it carries no servicing-income dimension.
  \item[Household debt histories] \citep{BartscherKuhnSchularickSteins2020} are the household
  sector alone and are distributional rather than servicing-income.
\end{description}

The claim is therefore ``none located under a stated protocol'', which is not ``none exists'';
Section~\ref{sec:limitations} gives the bounds of the search. What is claimed is accounting rather
than prediction: a pre-registered test of whether the accounts forecast bank-level credit losses
returned a null, and it is reported rather than left out. One priority claim in this line of work
was retired after a search found prior work on it.
